%============================================================
% constraints.tex — from the constraints to the algorithm (one-page guide)
% Grow the second table as chapters are written.
%============================================================
\chapter{From Constraints to Algorithm}
\label{ch:constraints}

A contest problem tells you its intended complexity through its limits. A judge runs roughly $10^8$ simple operations per second, so before choosing an algorithm, read $n$ (vertices), $m$ (edges) and $q$ (queries) and look the budget up here.

\section*{Step 1 --- What the limits allow}

\begingroup\small\renewcommand{\arraystretch}{1.3}
\begin{tabularx}{\linewidth}{@{}l l X@{}}
\toprule
\textbf{Limit} & \textbf{Budget} & \textbf{Typical graph algorithms} \\
\midrule
$n \le 10$--$12$ & $\bigO{n!}$, $\bigO{n^n}$ & try every order / every assignment; backtracking over paths \\
$n \le 20$ & $\bigO{2^n n^2}$ & Hamiltonian paths and TSP by bitmask DP, chromatic number by inclusion--exclusion (Ch.~\ref{ch:hard_small}) \\
$n \le 40$ & $\bigO{2^{n/2}}$ & meet in the middle: maximum clique / independent set (Ch.~\ref{ch:hard_small}) \\
$n \le 400$--$500$ & $\bigO{n^3}$ & Floyd--Warshall, transitive closure, all-pairs anything (Ch.~\ref{ch:all_pairs}); Hungarian algorithm (Ch.~\ref{ch:min_cost_flow}) \\
$n \le 5000$ & $\bigO{n^2}$ & adjacency matrix; BFS from every vertex; dense Prim / Dijkstra (Ch.~\ref{ch:dijkstra}, \ref{ch:mst}) \\
$n, m \le 10^5$--$2\cdot10^5$ & $\bigO{(n+m)\log n}$ & Dijkstra, Kruskal, LCA, HLD, small-to-large; $\bigO{(n+m)\log^2 n}$ usually still fits \\
$n, m \le 10^6$ & $\bigO{n+m}$ & BFS, iterative DFS, topological sort, DSU, Tarjan / Kosaraju --- watch recursion depth and memory \\
$n \le 10^{9}$ or more & $\bigO{\log n}$ per step & the graph is implicit: binary lifting (Ch.~\ref{ch:binary_lifting}), matrix powers (Ch.~\ref{ch:algebraic}), or a formula --- never store it \\
\bottomrule
\end{tabularx}
\endgroup

\begin{insightbox}[Read the graph's size, not only $n$]
An implicit graph's size is the number of \emph{states}: a $1000\times1000$ grid has $10^6$ vertices and about $4\cdot10^6$ edges; a lock with 4 wheels has $10^4$ states; ``position $\times$ keys held'' with 6 keys multiplies the vertices by $2^6 = 64$. Estimate $|V|$ and $|E|$ first (Chapter~\ref{ch:introduction}, Lesson~\ref{ssec:ch1-modelling}); the table above then applies to them.
\end{insightbox}

\clearpage
\section*{Step 2 --- What the question asks}

\begingroup\small\renewcommand{\arraystretch}{1.3}
\begin{tabularx}{\linewidth}{@{}>{\raggedright\arraybackslash}p{6.2cm} X@{}}
\toprule
\textbf{The problem says \ldots} & \textbf{Try first} \\
\midrule
``who points to whom'', ``how many connections'' & degrees only, no traversal (Ch.~\ref{ch:introduction}, P1) \\
``can $a$ reach $b$'', ``which rooms can be opened'' & DFS or BFS from the source (Ch.~\ref{ch:introduction}, P2) \\
grid of cells, ``regions'', ``islands'' & flood fill on the grid graph (Ch.~\ref{ch:introduction}, P3) \\
``minimum number of operations / moves / presses'' & BFS over states (Ch.~\ref{ch:introduction}, P4; Ch.~\ref{ch:bfs}) \\
``each element points to exactly one other'' & functional graph: follow arrows, find cycles (Ch.~\ref{ch:introduction}, P5); many jump queries: binary lifting (Ch.~\ref{ch:binary_lifting}) \\
cycles, subtrees, entry/exit times & DFS (Ch.~\ref{ch:dfs}) \\
unweighted shortest path, ``fewest steps'' & BFS (Ch.~\ref{ch:bfs}); weights 0 and 1: 0-1 BFS (Ch.~\ref{ch:negative_weights}) \\
``groups'', ``connected parts'', ``split into two teams'' & components, bipartite check (Ch.~\ref{ch:components}) \\
``prerequisites'', ``before / after'', ``order of tasks'' & topological sort, DP on a DAG (Ch.~\ref{ch:dag}) \\
``merge'', ``same group?'' online, edges only added & DSU (Ch.~\ref{ch:dsu}) \\
non-negative weights, ``cheapest / fastest route'' & Dijkstra (Ch.~\ref{ch:dijkstra}) \\
negative weights, ``arbitrage'', ``infinite profit'' & Bellman--Ford, negative cycles (Ch.~\ref{ch:negative_weights}) \\
``connect all at minimum cost'' & minimum spanning tree (Ch.~\ref{ch:mst}) \\
a tree with path / subtree queries & LCA, Euler tour, HLD, centroid decomposition (Ch.~\ref{ch:tree_basics}--\ref{ch:centroid}) \\
``critical road / server'', ``single point of failure'' & bridges and articulation points (Ch.~\ref{ch:bridges}) \\
``mutually reachable'', boolean constraints on pairs & SCC, 2-SAT (Ch.~\ref{ch:scc}--\ref{ch:two_sat}) \\
``use every edge exactly once'' & Euler path (Ch.~\ref{ch:euler}) \\
``assign / pair up'', ``maximum number of disjoint \ldots'' & matching and flows (Ch.~\ref{ch:max_flow}--\ref{ch:general_matching}) \\
edges added and removed over time, offline queries & DSU with rollback, segment tree over time (Ch.~\ref{ch:offline_connectivity}) \\
\bottomrule
\end{tabularx}
\endgroup

\begin{warningbox}[Limits are hints, not proofs]
The table gives the \emph{intended} complexity in most problems, not a guarantee. Small $n$ with huge values may hide a different graph; a large $n$ with ``sum of $n$ over test cases $\le 2\cdot10^5$'' allows per-test linear work but not per-test $\bigO{n^2}$. Always estimate with the worst test case.
\end{warningbox}
